\section{Risultati sperimentali}
\label{sec:risultati}

Tutti e 6 gli scenari (2 mappe $\times$ 3 configurazioni) sono stati eseguiti in modalità headless con un limite massimo di 500 tick. In tutti i casi tutti e 10 gli oggetti sono stati consegnati entro il limite (recovery 100\%), quindi il confronto si concentra su efficienza temporale ed energetica.

\subsection{Le mappe di test}

Le due istanze di mappa rappresentano ambienti con complessità crescente.

\begin{figure}[H]
  \centering
  \begin{subfigure}[b]{0.47\textwidth}
    \centering
    \includegraphics[width=\textwidth]{../png_grids/A.png}
    \caption{Mappa A -- ostacoli radi, corridoi ampi}
    \label{fig:map_a}
  \end{subfigure}
  \hfill
  \begin{subfigure}[b]{0.47\textwidth}
    \centering
    \includegraphics[width=\textwidth]{../png_grids/B.png}
    \caption{Mappa B -- ostacoli densi, percorsi vincolati}
    \label{fig:map_b}
  \end{subfigure}
  \caption{Visualizzazione delle due istanze di mappa. I quadrati colorati ai bordi sono i magazzini; i punti dorati sono gli oggetti da raccogliere.}
  \label{fig:maps}
\end{figure}

\textbf{Mappa A} presenta una struttura relativamente aperta, con ostacoli distribuiti in modo da non creare colli di bottiglia significativi. Gli agenti possono muoversi quasi liberamente tra i quadranti. \textbf{Mappa B} ha invece una densità di muri notevolmente superiore, con corridoi stretti che costringono gli agenti a percorsi obbligati e rendono più difficile sia l'esplorazione sia l'incontro tra agenti per la comunicazione.

\subsection{Metriche numeriche}

\begin{table}[H]
  \centering
  \caption{Risultati sperimentali -- Mappa A}
  \label{tab:results_a}
  \begin{tabular}{lcccc}
    \toprule
    \textbf{Configurazione} & \textbf{Tick totali} & \textbf{Energia media} & \textbf{Energia/oggetto} & \textbf{Tick/oggetto} \\
    \midrule
    Exploration & 251 & 218.4 & 109.2 & 25.1 \\
    Collection  & 129 & 129.0 &  64.5 & 12.9 \\
    With Relay  & 133 & 133.0 &  66.5 & 13.3 \\
    \bottomrule
  \end{tabular}
\end{table}

\begin{table}[H]
  \centering
  \caption{Risultati sperimentali -- Mappa B}
  \label{tab:results_b}
  \begin{tabular}{lcccc}
    \toprule
    \textbf{Configurazione} & \textbf{Tick totali} & \textbf{Energia media} & \textbf{Energia/oggetto} & \textbf{Tick/oggetto} \\
    \midrule
    Exploration & 153 & 152.6 & 76.3 & 15.3 \\
    Collection  & 117 & 117.0 & 58.5 & 11.7 \\
    With Relay  & 117 & 117.0 & 58.5 & 11.7 \\
    \bottomrule
  \end{tabular}
\end{table}

\paragraph{Lettura dei dati.}
La configurazione \emph{Exploration} è sistematicamente la peggiore in entrambe le mappe: 251 tick su A e 153 su B, con energia per oggetto quasi doppia rispetto alle altre. Il surplus di Scout non compensa la carenza di Collector nella fase operativa.

\emph{Collection} e \emph{With Relay} sono invece molto competitive tra loro. Su Mappa~A, Collection è marginalmente migliore (129 vs.\ 133 tick, differenza del 3\%); su Mappa~B le due configurazioni pareggiano esattamente. Questo risultato merita un'analisi più approfondita (si veda la Sezione~\ref{sec:relayvsmap}).

\subsection{Grafici comparativi}

\subsubsection{Consegne cumulative nel tempo}

\begin{figure}[H]
  \centering
  \includegraphics[width=0.9\textwidth]{../experiments/comparison_objects_20260318_160902.png}
  \caption{Numero cumulativo di oggetti consegnati nel tempo, per tutte le configurazioni e mappe. Le curve più ripide indicano un ritmo di consegna più sostenuto.}
  \label{fig:comparison_objects}
\end{figure}

Il grafico mostra chiaramente come la configurazione \emph{Exploration} (curve con pendenza minore) accumuli ritardi già nelle prime fasi: la scoperta degli oggetti è rapida, ma la consegna è lenta per mancanza di Collector. Al contrario, \emph{Collection} e \emph{With Relay} mantengono una pendenza più ripida e costante, indicando un ritmo di consegna regolare.

\begin{figure}[H]
  \centering
  \includegraphics[width=0.9\textwidth]{../experiments/subplots_objects_20260318_160902.png}
  \caption{Griglia di subplot $2 \times 3$: righe = istanze (A, B), colonne = configurazioni (Exploration, Collection, With Relay). Ogni pannello mostra la curva cumulativa degli oggetti consegnati per quella combinazione.}
  \label{fig:subplots_objects}
\end{figure}

Osservando i subplot per istanza, emerge un pattern interessante su Mappa~B: la configurazione \emph{With Relay} raggiunge il completamento ($10/10$ oggetti) con la stessa traiettoria di \emph{Collection}, ma con una distribuzione temporale delle consegne leggermente più uniforme nelle fasi centrali, suggerendo un flusso informativo più regolare.

\subsubsection{Consumo energetico nel tempo}

\begin{figure}[H]
  \centering
  \includegraphics[width=0.9\textwidth]{../experiments/comparison_energy_20260318_160902.png}
  \caption{Energia totale consumata dal sistema nel tempo (somma su tutti gli agenti), per tutte le configurazioni e mappe.}
  \label{fig:comparison_energy}
\end{figure}

\begin{figure}[H]
  \centering
  \includegraphics[width=0.9\textwidth]{../experiments/subplots_energy_20260318_160902.png}
  \caption{Griglia di subplot per il consumo energetico: ogni pannello mostra il profilo di consumo di quella combinazione istanza-configurazione.}
  \label{fig:subplots_energy}
\end{figure}

Il profilo energetico è lineare per costruzione (ogni agente consuma 1 unità per tick finché è attivo). L'energia totale consumata a fine simulazione è quindi proporzionale al numero di tick. Il grafico evidenzia come le curve di \emph{Exploration} abbiano una pendenza moderata ma si estendano molto più a lungo, mentre \emph{Collection} e \emph{With Relay} terminano bruscamente a energie totali molto più basse.

\subsection{Heatmap di visita}

Le heatmap mostrano la frequenza di visita di ciascuna cella della griglia aggregata su tutti gli agenti e tutti i tick della simulazione. Zone più calde (rosse) indicano celle frequentate spesso, zone fredde (blu/viola) celle raramente visitate.

\subsubsection{Mappa A}

\begin{figure}[H]
  \centering
  \begin{subfigure}[b]{0.32\textwidth}
    \includegraphics[width=\textwidth]{../experiments/A/20260318_160902_exploration/heatmap.png}
    \caption{Exploration}
  \end{subfigure}
  \hfill
  \begin{subfigure}[b]{0.32\textwidth}
    \includegraphics[width=\textwidth]{../experiments/A/20260318_160903_collection/heatmap.png}
    \caption{Collection}
  \end{subfigure}
  \hfill
  \begin{subfigure}[b]{0.32\textwidth}
    \includegraphics[width=\textwidth]{../experiments/A/20260318_160903_with_relay/heatmap.png}
    \caption{With Relay}
  \end{subfigure}
  \caption{Heatmap di visita per Mappa~A nelle tre configurazioni.}
  \label{fig:heatmap_a}
\end{figure}

Su Mappa~A, la configurazione \emph{Exploration} produce heatmap con una distribuzione più uniforme della presenza agenti sull'intera griglia: gli Scout diffondono sistematicamente la copertura su tutta la mappa. La configurazione \emph{Collection} mostra invece una concentrazione più elevata nelle aree prossime ai magazzini e agli oggetti, con meno ridondanza nelle zone già esplorate. La configurazione \emph{With Relay} evidenzia una \emph{striscia} di attività nel centro della mappa, traccia visibile del percorso a diamante del Relay.

\subsubsection{Mappa B}

\begin{figure}[H]
  \centering
  \begin{subfigure}[b]{0.32\textwidth}
    \includegraphics[width=\textwidth]{../experiments/B/20260318_160904_exploration/heatmap.png}
    \caption{Exploration}
  \end{subfigure}
  \hfill
  \begin{subfigure}[b]{0.32\textwidth}
    \includegraphics[width=\textwidth]{../experiments/B/20260318_160904_collection/heatmap.png}
    \caption{Collection}
  \end{subfigure}
  \hfill
  \begin{subfigure}[b]{0.32\textwidth}
    \includegraphics[width=\textwidth]{../experiments/B/20260318_160905_with_relay/heatmap.png}
    \caption{With Relay}
  \end{subfigure}
  \caption{Heatmap di visita per Mappa~B nelle tre configurazioni.}
  \label{fig:heatmap_b}
\end{figure}

Su Mappa~B l'effetto della densità degli ostacoli è evidente: le heatmap mostrano \textbf{corridoi obbligati} di transito ad alta frequenza, separati da zone completamente fredde (i muri). La configurazione \emph{With Relay} produce una distribuzione più bilanciata dell'attività lungo i corridoi centrali, a conferma del ruolo di ponte svolto dall'agente Relay tra le zone periferiche degli Scout e le aree magazzino dei Collector.

\subsection{Relay vs.\ struttura della mappa}
\label{sec:relayvsmap}

I dati mostrano un pattern importante: \textbf{la superiorità del Relay cresce con la complessità topologica della mappa}.

Su \textbf{Mappa~A} (aperta), il Relay introduce un leggero overhead rispetto a \emph{Collection}: 133 tick vs.\ 129 tick (+3\%). In una mappa aperta, gli agenti si incontrano naturalmente e spesso senza bisogno di un intermediario dedicato: i Collector con raggio di comunicazione 1 già ricevono informazioni dagli Scout nelle zone centrali con sufficiente frequenza. Il Relay diventa quasi ridondante, e il costo di avere un Collector in meno si fa sentire.

Su \textbf{Mappa~B} (vincolata), il Relay recupera completamente questo deficit: le due configurazioni pareggiano a 117 tick. La spiegazione è nei corridoi stretti: in un ambiente labirintico, Scout e Collector raramente si incontrano direttamente nelle traiettorie ottimali, e il Relay -- che pattuglia il centro, punto di passaggio obbligato -- risolve questo collo di bottiglia comunicativo. La perdita di un Collector è bilanciata esattamente dalla maggiore velocità di propagazione delle informazioni sugli oggetti.

L'ipotesi naturale è che su mappe con complessità ancora maggiore (più ostacoli, corridoi più stretti, maggior distanza tra zone esplorazione e magazzini), il vantaggio del Relay diventerebbe più marcato e la configurazione \emph{With Relay} supererebbe chiaramente \emph{Collection}.
